\chapter{Trabalhos relacionados}
\label{textuais/cap2}
Este capítulo apresenta soluções existentes voltadas ao monitoramento e à
automação da irrigação, com o objetivo de identificar suas características,
possibilidades e limitações, e de situar a proposta deste trabalho em relação a
elas. Foram considerados tanto produtos disponíveis no mercado quanto sistemas
desenvolvidos em pesquisas acadêmicas, uma vez que os primeiros mostram o que
já está acessível ao produtor, enquanto os segundos indicam as abordagens
técnicas mais recentes.

Cada solução é descrita em uma seção própria. As Seções~\ref{sec:hydrawise},
\ref{sec:netbeat} e~\ref{sec:valley} apresentam três soluções comerciais, e as
Seções~\ref{sec:campos}, \ref{sec:pereira} e~\ref{sec:abdelmoneim} apresentam três trabalhos acadêmicos de baixo custo. Por
fim, a Seção~\ref{sec:comparativo-trabalhos} compara essas soluções.

% =====================================================================

\section{Hydrawise}
\label{sec:hydrawise}

O Hydrawise é uma solução comercial da Hunter Industries: um software de
gerenciamento de irrigação que funciona em conjunto com controladores
conectados por Wi-Fi da própria fabricante. Seu principal recurso, chamado
Predictive Watering, ajusta automaticamente os programas de irrigação
com base na previsão do tempo local, considerando temperatura, probabilidade de
chuva, velocidade do vento e umidade do ar. Esses dados vêm de estações
meteorológicas virtuais, que combinam informações de satélites, estações
físicas e outras fontes \cite{hunterhydrawise}.

O sistema pode ser operado remotamente por aplicativo ou navegador, permite o
acionamento manual de setores e emite alertas como os de vazamento. Segundo o
fabricante, o ajuste preditivo pode economizar até 50\% de água. A solução é
voltada principalmente a jardins residenciais e áreas paisagísticas
comerciais, com um plano gratuito básico e planos pagos que ampliam o histórico
de dados e o número de controladores \cite{hunterhydrawise}.

% =====================================================================
\section{NetBeat}
\label{sec:netbeat}

O NetBeat é uma solução comercial da Netafim, apresentada em 2018 como uma
plataforma que integra monitoramento, análise e automação da irrigação e da
fertirrigação. O sistema reúne dados de sensores de solo, de planta e
hidráulicos, além de informações climáticas, e os processa em nuvem por meio de
modelos dinâmicos de cultura desenvolvidos pela empresa, que recomendam quando
irrigar e fertirrigar ao longo do ciclo da cultura
\cite{israelagri2018netbeat}.

Sua arquitetura é composta por uma unidade de controle principal e por unidades
terminais remotas, que se comunicam por diferentes tecnologias, como Wi-Fi,
redes celulares e LoRa, e podem ser alimentadas por energia solar. O controle
pode ser feito no local ou remotamente, por computador ou celular
\cite{e4cnetbeat}. Atualmente, a página do produto apresenta a plataforma
GrowSphere como a solução de agricultura digital da empresa
\cite{netafimgrowsphere}.

% =====================================================================
\section{Valley Scheduling}
\label{sec:valley}

O Valley Scheduling é uma solução comercial da Valmont Industries: um software
de programação da irrigação lançado inicialmente no mercado brasileiro e
expandido para a América do Norte em 2018. Ele gera recomendações a partir do
tipo de solo, da cultura, da fase de desenvolvimento e de dados meteorológicos,
sendo compatível com pivôs centrais, sistemas lineares e gotejamento. Na época
da expansão, a ferramenta já havia sido utilizada em cerca de 5 milhões de
acres em 12 países \cite{valmont2018}.

As recomendações podem ser modeladas, por meio de um algoritmo proprietário que
dispensa sensores no campo, ou medidas, com dados de sondas de umidade do solo e
estações meteorológicas. O produtor recebe a necessidade hídrica da cultura e
uma previsão de irrigação para cinco a sete dias, visualizada em listas ou mapas
pelo celular ou computador, com acompanhamento de um especialista em agronomia
\cite{valleyscheduling}.

% =====================================================================
\section{Plataforma de baixo custo e código aberto para automação da irrigação}
\label{sec:campos}

Diferentemente das soluções anteriores, trata-se de um trabalho acadêmico.
\citeonline{campos2021} desenvolveram, no Instituto Federal Goiano - Campus Ceres, uma
plataforma de baixo custo e código aberto para automatizar a irrigação com base
na umidade do solo. O sistema utilizou um Arduino Mega 2560 conectado a quatro
sensores capacitivos de umidade do solo (10HS), a um sensor de temperatura e
umidade do ar (DHT22), a um display LCD, a um módulo de relógio de
tempo real, a um cartão de memória e a um módulo \sigla{GSM}{Global System for Mobile Communications} para envio de mensagens
\sigla{SMS}{Short Message Service}. Por meio de relés, o controlador acionava válvulas solenoides sempre que a
umidade do solo ficava abaixo de um limiar predefinido.

O sistema foi avaliado no cultivo de tomate-cereja em casa de vegetação, com
quatro limiares de umidade volumétrica (0,23 a 0,44~m$^3$~m$^{-3}$). Após a
calibração para o solo, as leituras dos sensores alcançaram índice de desempenho
classificado como ótimo ($c = 0{,}95$), com desvio absoluto médio de 1,34\% em
relação ao método gravimétrico. O aumento do limiar de 0,23 para
0,44~m$^3$~m$^{-3}$ elevou em 50,52\% a eficiência do uso da água
\cite{campos2021}.

Os autores também apontam limitações. O custo total foi de R\$~1.390,55, dos
quais R\$~978,35 corresponderam apenas aos quatro sensores de umidade. Os dados
precisavam ser recuperados manualmente do cartão de memória, o que consome tempo
e aumenta o risco de perdas, e os autores recomendam o uso de Wi-Fi para enviar
os dados à nuvem e acompanhá-los em tempo real. Por fim, a plataforma exige mão
de obra especializada para desenvolvimento e manutenção \cite{campos2021}.

% =====================================================================
\section{Sistema de gotejamento inteligente com ESP32}
\label{sec:pereira}

\citeonline{pereira2023} desenvolveram um sistema de irrigação por gotejamento baseado no ESP32, escolhido pelo baixo custo e pelo Wi-Fi integrado. O dispositivo usa um sensor capacitivo de umidade do solo, um sensor de temperatura do solo, um sensor DHT22 de temperatura e umidade do ar e um sensor de vazão, e aciona por relé uma válvula solenoide e uma bomba. Os dados são enviados por Wi-Fi à plataforma Blynk, cujo aplicativo e painel web exibem gráficos das leituras, guardam os registros de cada irrigação e permitem abrir a válvula manualmente ou desligar o modo automático.

A decisão de irrigar é feita por limiar. A válvula é aberta quando o solo está seco, dentro de janelas de horário pela manhã e à noite, e fica aberta por um tempo fixo de uma hora. O sistema foi avaliado por uma semana no cultivo de cebolinha em um pequeno recipiente. Os próprios autores relatam excesso de irrigação causado pela demora do sensor em perceber a água aplicada, problema que levaram a corrigir no firmware, bloqueando uma segunda irrigação no mesmo dia. A economia de água não foi medida, e o custo do sistema não foi informado \cite{pereira2023}.

% =====================================================================
\section{Manejo por umidade do solo e por clima na alface}
\label{sec:abdelmoneim}

\citeonline{abdelmoneim2025comparative} construíram um dispositivo de baixo custo, com um ESP32 e quatro sensores capacitivos instalados em diferentes profundidades, alimentado por bateria e painel solar, que envia as leituras por Wi-Fi à plataforma ThingSpeak a cada seis horas. O protótipo custou cerca de 62 dólares, contra 800 a 1.500 dólares estimados para equipamentos comerciais equivalentes. Os dados podem ser consultados em gráficos na plataforma e exportados em arquivos.

O dispositivo foi usado para comparar duas formas de manejo da irrigação por gotejamento de alface em campo, na Itália. Em uma parcela, a irrigação era iniciada quando a umidade do solo atingia um limiar. Na outra, seguia a programação pelo método FAO-56, com a ET$_0$ de Penman-Monteith e coeficientes de cultura. O manejo pela umidade do solo usou 28,8\% menos água e obteve produtividade da água 52,5\% maior. Os autores apontam como limitações a avaliação em uma única safra e um único tipo de solo, sem repetições, e a calibração dos sensores específica para aquele solo \cite{abdelmoneim2025comparative}.

% =====================================================================
\section{Comparativo entre os trabalhos}
\label{sec:comparativo-trabalhos}

A Tabela~\ref{tab:comparativo-trabalhos} compara as soluções apresentadas a
partir de oito critérios. Cada critério é explicado a seguir.

\begin{table}[htbp]
  \centering
  \caption{Comparativo entre os trabalhos relacionados}
  \label{tab:comparativo-trabalhos}
  \scriptsize
  \renewcommand{\arraystretch}{1.25}
  \setlength{\tabcolsep}{3pt}
  \begin{tabularx}{\textwidth}{>{\raggedright\arraybackslash}p{1.6cm}
                                >{\raggedright\arraybackslash}X
                                >{\raggedright\arraybackslash}X
                                >{\raggedright\arraybackslash}X
                                >{\raggedright\arraybackslash}X
                                >{\raggedright\arraybackslash}X
                                >{\raggedright\arraybackslash}X}
    \hline
    \textbf{Critério} & \textbf{Hydrawise} & \textbf{NetBeat} &
    \textbf{Valley Scheduling} & \textbf{\citeonline{campos2021}} &
    \textbf{\citeonline{pereira2023}} & \textbf{\citeonline{abdelmoneim2025comparative}}\\
    \hline
    Público-alvo &
    Residências e profissionais de paisagismo &
    Desde grandes áreas até pequenos produtores &
    Produtores rurais com pivô, linear ou gotejamento &
    Pesquisa em casa de vegetação &
    Pequenos cultivos, avaliado em recipiente &
    Pesquisa em campo, com alface \\
    \hline
    Base para decidir a irrigação &
    Possibilita ajustar os programas pela previsão do tempo &
    Possibilita recomendações por modelos dinâmicos de cultura &
    Possibilita recomendações por algoritmo ou por dados medidos &
    Aciona quando a umidade do solo cai abaixo do limiar &
    Limiar de umidade do solo, com janelas de horário &
    Compara limiar de umidade do solo e FAO-56 \\
    \hline
    Variáveis consideradas &
    Dados meteorológicos externos; pode receber sensores &
    Solo, planta, clima e hidráulica &
    Solo, cultura, fase e clima &
    Umidade do solo, temperatura e umidade do ar &
    Umidade e temperatura do solo, temperatura e umidade do ar e vazão &
    Umidade do solo em quatro profundidades \\
    \hline
    Forma de atuação &
    Possibilita controle automático e manual dos setores &
    Possibilita controlar irrigação e fertirrigação &
    Fornece recomendação; a aplicação fica a cargo do produtor &
    Acionamento de válvulas por relés &
    Válvula e bomba por relé, automático e manual &
    Apenas monitoramento \\
    \hline
    Acesso remoto &
    Aplicativo e navegador &
    Computador e celular &
    Celular, tablet e computador &
    Apenas SMS com o estado da irrigação &
    Aplicativo e painel web (Blynk) &
    Painel web (ThingSpeak) \\
    \hline
    Histórico de dados &
    Varia conforme o plano contratado &
    Dados processados em nuvem &
    Relatório ao fim da safra &
    Cartão de memória, com coleta manual &
    Gráficos e registros de irrigação no Blynk &
    Gráficos e exportação de arquivos \\
    \hline
    Comunicação &
    Wi-Fi &
    Wi-Fi, celular, LoRa e outras &
    Internet, conforme os equipamentos integrados &
    GSM (SMS) &
    Wi-Fi &
    Wi-Fi, a cada seis horas \\
    \hline
    Abertura da solução &
    Proprietária &
    Proprietária &
    Proprietária &
    Código aberto &
    Não informada, usa a plataforma Blynk &
    Protótipo próprio, de cerca de US\$~62 \\
    \hline
  \end{tabularx}
  \fonte{Elaborado pelo autor com base em \citeonline{hunterhydrawise},
  \citeonline{israelagri2018netbeat}, \citeonline{e4cnetbeat},
  \citeonline{valmont2018}, \citeonline{valleyscheduling},
  \citeonline{campos2021}, \citeonline{pereira2023} e \citeonline{abdelmoneim2025comparative}}
\end{table}

O público-alvo indica a quem cada solução se destina e em que escala
ela opera. O Hydrawise atende residências e profissionais de paisagismo, o
NetBeat abrange desde grandes áreas até pequenos produtores, e o Valley
Scheduling é voltado a produtores rurais que irrigam por pivô, sistema linear ou
gotejamento. A plataforma de \citeonline{campos2021} foi concebida para pesquisa
em casa de vegetação e estufas. O sistema de \citeonline{pereira2023} foi pensado para pequenos cultivos e avaliado em um recipiente, e o dispositivo de \citeonline{abdelmoneim2025comparative} foi usado em uma pesquisa de campo com alface.

A base para decidir a irrigação corresponde ao critério usado para
definir quando e quanto irrigar. O Hydrawise ajusta os programas pela previsão
do tempo, o NetBeat recomenda a irrigação por modelos dinâmicos de cultura e o
Valley Scheduling gera recomendações por algoritmo ou por dados medidos em
campo. A plataforma de \citeonline{campos2021} apenas aciona a irrigação quando
a umidade do solo cai abaixo de um limiar, e o sistema de \citeonline{pereira2023} faz o mesmo, mas apenas em janelas de horário e por um tempo fixo. Já \citeonline{abdelmoneim2025comparative} comparam as duas abordagens, a do limiar de umidade e a da FAO-56, sem combiná-las.

As variáveis consideradas são as informações que alimentam essa
decisão. As soluções comerciais utilizam um conjunto mais amplo de dados, que
inclui informações de solo, de cultura, de clima e, no caso do NetBeat, de
sensores de planta e hidráulicos. A plataforma de \citeonline{campos2021} mede apenas três variáveis, a umidade do
solo, a temperatura e a umidade do ar. O sistema de \citeonline{pereira2023} acrescenta a temperatura do solo e a vazão, e o dispositivo de \citeonline{abdelmoneim2025comparative} mede apenas a umidade do solo, em quatro profundidades.

A forma de atuação indica se a solução executa a irrigação ou apenas
orienta o produtor. O Hydrawise e o NetBeat controlam diretamente os setores de
irrigação, e o NetBeat também a fertirrigação. O Valley Scheduling apenas
fornece recomendações, cabendo ao produtor aplicá-las. A plataforma de
\citeonline{campos2021} aciona válvulas por meio de relés, e o sistema de \citeonline{pereira2023} aciona uma válvula e uma bomba, nos modos automático e manual. O dispositivo de \citeonline{abdelmoneim2025comparative} apenas monitora o solo.

O acesso remoto refere-se à possibilidade de acompanhar e comandar o
sistema a distância. As três soluções comerciais oferecem acesso por aplicativo,
navegador ou computador. A plataforma de \citeonline{campos2021} limita-se ao
envio de mensagens SMS com o estado da irrigação, sem permitir comandos. O sistema de \citeonline{pereira2023} oferece aplicativo e painel web pela plataforma Blynk, que também permitem comandar a válvula, e o de \citeonline{abdelmoneim2025comparative} permite apenas consultar os dados na plataforma ThingSpeak.

O histórico de dados indica como as informações ficam disponíveis para
consulta posterior. No Hydrawise, a extensão do histórico depende do plano
contratado; no NetBeat, os dados são processados em nuvem; e no Valley
Scheduling, o produtor recebe um relatório ao fim da safra. Na plataforma de
\citeonline{campos2021}, os dados ficam em um cartão de memória e precisam ser
coletados manualmente. Nos trabalhos de \citeonline{pereira2023} e de \citeonline{abdelmoneim2025comparative}, o histórico fica nas plataformas Blynk e ThingSpeak, respectivamente.

A comunicação descreve a tecnologia usada para transmitir os dados. O
Hydrawise utiliza Wi-Fi; o NetBeat combina Wi-Fi, redes celulares, LoRa e outras
tecnologias; e o Valley Scheduling depende da conexão à internet dos
equipamentos integrados. A plataforma de \citeonline{campos2021} usa a rede GSM
para enviar SMS, e os trabalhos de \citeonline{pereira2023} e de \citeonline{abdelmoneim2025comparative} usam o Wi-Fi, este último enviando os dados apenas a cada seis horas para economizar energia.

Por fim, a abertura da solução indica se o código e o projeto estão
disponíveis para adaptação. As três soluções comerciais são proprietárias,
enquanto a plataforma de \citeonline{campos2021} é de código aberto, o que
permite substituir sensores e ajustar a lógica de
irrigação conforme a necessidade. Os trabalhos de \citeonline{pereira2023} e de \citeonline{abdelmoneim2025comparative} não informam se o código está disponível, e o primeiro depende de uma plataforma de terceiros para a interface.

Em síntese, as soluções comerciais oferecem modelos de decisão mais elaborados
e suporte técnico, mas são proprietárias e voltadas a escalas específicas. Os trabalhos acadêmicos mostram a viabilidade de soluções de baixo custo com hardware aberto e, no caso de \citeonline{pereira2023} e \citeonline{abdelmoneim2025comparative}, com o ESP32 e o envio dos dados pela internet. Nenhum deles, porém, combina o limiar de umidade, para decidir quando irrigar, com o cálculo pela FAO-56, para decidir quanto irrigar, e apenas o de \citeonline{pereira2023} permite comandar a irrigação a distância, por meio de uma plataforma de terceiros.